\section{MI-04: a characterization of essentially Hermitian off-diagonal blocks}
\label{sec:mi04}

A matrix $X\in\M_n(\C)$ is \emph{essentially Hermitian} if
$X=\alpha H+\beta I$ for some Hermitian $H$ and complex scalars $\alpha,\beta$.
The target is the necessity assertion in the following theorem, corresponding to Bourin and Lee's Conjecture 3.3 \cite{BL2021}.

\begin{theorem}\label{thm:mi04}
For $X\in\M_n(\C)$, the following conditions are equivalent.
\begin{enumerate}
\item For every positive semidefinite completion
$P=\begin{pmatrix}A&X\\X^*&B\end{pmatrix}$,
\[
 \opnorm{P}\leq\opnorm{A+B}.
\]
\item The same inequality holds for every such completion whose diagonal-block sum $A+B$ is a scalar multiple of $I$.
\item For every Hermitian $T$, the largest and smallest eigenvalues of
\[
 K_X(T)=\begin{pmatrix}T&X\\X^*&-T\end{pmatrix}
\]
are negatives of one another.
\item $X$ is essentially Hermitian.
\end{enumerate}
In particular, no invertibility or distinct-singular-value hypothesis on $X$ is required.
\end{theorem}

\subsection{Scalar-sum completions force extreme spectral symmetry}

The implication (1)$\Rightarrow$(2) is immediate. Assume (2). Fix a Hermitian $T$ and write
\[
 a=\lambda_{\max}(K_X(T)),\qquad b=-\lambda_{\min}(K_X(T)).
\]
Since $\tr K_X(T)=0$, both $a$ and $b$ are nonnegative. The matrix
\[
 bI_{2n}+K_X(T)
 =\begin{pmatrix}bI+T&X\\X^*&bI-T\end{pmatrix}
\]
is positive semidefinite, has diagonal-block sum $2bI$, and has norm $a+b$. Thus (2) gives $a\leq b$.

Let $J=\diag(I,-I)$. Direct multiplication shows
\[
 K_X(-T)=J[-K_X(T)]J^*.
\]
Applying the same argument to $-T$ gives $b\leq a$. This proves (3).
If ``positive'' is interpreted as positive definite, use $b+\eta$ in place of $b$ and let $\eta\downarrow0$; the same conclusion follows.

\subsection{A second-order test forces symmetry of off-diagonal magnitudes}

Assume (3). It continues to hold with $X$ replaced by $\varepsilon X$, for every $\varepsilon>0$, because
\[
 K_{\varepsilon X}(T)=\varepsilon K_X(T/\varepsilon).
\]
It is also invariant under unitary similarity of $X$.

Choose any orthonormal basis and write $x_{ij}$ for the entries of $X$ in that basis. Let
\[
 D=\diag(1,d_2,\ldots,d_n),\qquad -1<d_j<1,
\]
and consider
\[
 K(\varepsilon)=\begin{pmatrix}D&\varepsilon X\\\varepsilon X^*&-D\end{pmatrix}.
\]
At $\varepsilon=0$, its largest and smallest eigenvalues are the simple eigenvalues $1$ and $-1$. For small $\varepsilon$, denote their continuations by $\lambda_+(\varepsilon)$ and $\lambda_-(\varepsilon)$. The elementary second-order perturbation formula gives, with $d_1=1$,
\begin{align}
 \lambda_+(\varepsilon)
 &=1+\varepsilon^2\sum_{j=1}^n\frac{|x_{1j}|^2}{1+d_j}
   +O(\varepsilon^3),\label{eq:mi04-plus}\\
 \lambda_-(\varepsilon)
 &=-1-\varepsilon^2\sum_{j=1}^n\frac{|x_{j1}|^2}{1+d_j}
   +O(\varepsilon^3).\label{eq:mi04-minus}
\end{align}
For completeness, if $H_0u=\lambda_0u$ is a simple eigenpair of a Hermitian matrix, $\|u\|=1$, and $V$ is Hermitian, expand an eigenvector of $H_0+\varepsilon V$ as $u+\varepsilon w+O(\varepsilon^2)$ with $w\perp u$. In an orthonormal eigenbasis $(u_j)$ of $H_0$,
\[
 w=\sum_{u_j\perp u}\frac{u_j^*Vu}{\lambda_0-\lambda_j}u_j,
\]
and the coefficient of $\varepsilon^2$ in the eigenvalue is
$u^*Vw=\sum_{u_j\perp u}|u_j^*Vu|^2/(\lambda_0-\lambda_j)$.
Applying this to the off-diagonal perturbation gives \eqref{eq:mi04-plus}--\eqref{eq:mi04-minus}.

Condition (3) implies $\lambda_+(\varepsilon)+\lambda_-(\varepsilon)=0$. Comparing the quadratic coefficients and canceling the $j=1$ terms gives
\begin{equation}\label{eq:mi04-weights}
 \sum_{j=2}^n\frac{|x_{1j}|^2-|x_{j1}|^2}{1+d_j}=0
 \qquad\text{for every }(d_2,\ldots,d_n)\in(-1,1)^{n-1}.
\end{equation}
First take all $d_j=0$. Then change only $d_\ell$ to any nonzero number in $(-1,1)$. Subtracting the two identities proves
\[
 |x_{1\ell}|=|x_{\ell1}|\qquad(\ell\geq2).
\]
The basis was arbitrary, so
\begin{equation}\label{eq:mi04-pair}
 |u^*Xv|=|v^*Xu|
 \quad\text{for all orthonormal pairs }u,v.
\end{equation}

\subsection{The pair condition implies normality and a collinear spectrum}

Complete a unit vector $u$ to an orthonormal basis and sum the squares in \eqref{eq:mi04-pair}. Including the equal diagonal terms yields
\[
 \|X^*u\|^2=\|Xu\|^2.
\]
Thus $u^*(XX^*-X^*X)u=0$ for every $u$, and $X$ is normal.

Use an orthonormal eigenbasis of $X$, with eigenvalues $z_1,\ldots,z_n$. If $n\leq2$, these eigenvalues lie on a line. For any triple of indices $i,j,k$ when $n\geq3$, let $\omega=e^{2\pi i/3}$ and set
\[
 u=\frac{e_i+e_j+e_k}{\sqrt3},\qquad
 v=\frac{e_i+\omega e_j+\omega^2e_k}{\sqrt3}.
\]
These vectors are orthonormal. Equation \eqref{eq:mi04-pair} gives
\[
 |z_i+\omega z_j+\omega^2z_k|
 =|z_i+\omega^2z_j+\omega z_k|.
\]
Writing $a=z_i-z_k$ and $b=z_j-z_k$, the difference of the squared sides is
\[
 |a+\omega b|^2-|a+\omega^2b|^2
 =2\sqrt3\,\operatorname{Im}(a\overline b).
\]
It follows that every eigenvalue triple is collinear. If two eigenvalues are distinct, all the others lie on the line through those two; if none are distinct, $X$ is scalar. Since $X$ is normal, in either case $X=\alpha H+\beta I$ with $H$ Hermitian. This proves (3)$\Rightarrow$(4).

\subsection{Converse and completion of the equivalence}

Assume (4). Multiplying $X$ by a scalar of modulus one changes $K_X(T)$ only by unitary similarity. We may therefore write
\[
 X=H+ibI,\qquad H=H^*,\quad b\in\R.
\]
With the unitary matrix
\[
 Q=\frac1{\sqrt2}\begin{pmatrix}I&I\\ iI&-iI\end{pmatrix},
\]
direct multiplication gives
\[
 Q^*K_X(T)Q
 =\begin{pmatrix}-bI&T-iH\\T+iH&bI\end{pmatrix}.
\]
A singular value decomposition of $T-iH$ reduces this matrix to $n$ blocks of the form
$\begin{pmatrix}-b&\sigma_j\\\sigma_j&b\end{pmatrix}$.
Its eigenvalues are therefore the pairs
$\pm\sqrt{\sigma_j^2+b^2}$. In particular, (3) holds.

Finally assume (3) and let
$P=\begin{pmatrix}A&X\\X^*&B\end{pmatrix}\geq0$.
Set $c=\opnorm{A+B}$. Because $A+B\leq cI$, the matrix
\[
 P'=\begin{pmatrix}A&X\\X^*&cI-A\end{pmatrix}
 =P+\diag(0,cI-A-B)
\]
satisfies $P'\geq P\geq0$. Also
\[
 P'=\frac c2I_{2n}+K_X\left(A-\frac c2I\right).
\]
Positivity gives $\lambda_{\min}(K_X(A-cI/2))\geq-c/2$. By (3), its largest eigenvalue is at most $c/2$. Thus $P'\leq cI_{2n}$ and consequently $P\leq cI_{2n}$. This proves (1) and finishes the proof of Theorem~\ref{thm:mi04}.
\hfill$\square$
